% Part: applied-modal-logics
% Chapter: epistemic-logic
% Section: language-epistemic-logic

\documentclass[../../../include/open-logic-section]{subfiles}

\begin{document}

\olfileid{aml}{el}{pal}

\olsection{公开宣告逻辑}

动态认知逻辑使我们能够表示主体的知识如何随时间变化，或如何因获得新信息而改变。
其中许多逻辑使用信息\emph{事件}或\emph{更新}来表示知识的变化。
最基本的一类更新是\emph{公开宣告}，即某个公式被真实宣告，且所有主体共同见证这一过程。
为此，我们如下扩充语言：

\begin{defn}
设 $G$ 是一个主体符号的集合。带公开宣告的多主体认知逻辑的基本语言包含：
\begin{enumerate}
  \tagitem{prvFalse}{表示假的命题常项~$\lfalse$}{}
  \tagitem{prvTrue}{表示真的命题常项~$\ltrue$}{}
  \item 一个可数无穷的命题变元集合：$\Obj p_0$, $\Obj p_1$, $\Obj p_2$, \dots
  \item 命题联结词：\startycommalist
  \iftag{prvNot}{\ycomma $\lnot$（否定）}{}%
  \iftag{prvAnd}{\ycomma $\land$（合取）}{}%
  \iftag{prvOr}{\ycomma $\lor$（析取）}{}%
  \iftag{prvIf}{\ycomma $\lif$（条件句）}{}%
  \item 知识算子 $\Knows_a$，其中 $a \in G$。
  \item 公开宣告算子 $[!B]$，其中 $!B$ 是公式。
\end{enumerate}
\end{defn}

公开宣告算子的作用类似于方框算子，语言的归纳定义相应给出如下：

\begin{defn}
  认知语言的\emph{公式}归纳定义如下：
  \begin{enumerate}
  \tagitem{prvFalse}{$\lfalse$ 是原子公式。}{}

  \tagitem{prvTrue}{$\ltrue$ 是原子公式。}{}

  \item 每个命题变元 $\Obj p_i$ 是（原子）公式。

  \tagitem{prvNot}{若 $!A$ 是公式，则 $\lnot !A$ 是公式。}{}

  \tagitem{prvAnd}{若 $!A$ 和 $!B$ 是公式，则 $(!A \land !B)$ 是公式。}{}

  \tagitem{prvOr}{若 $!A$ 和 $!B$ 是公式，则 $(!A \lor !B)$ 是公式。}{}

  \tagitem{prvIf}{若 $!A$ 和 $!B$ 是公式，则 $(!A \lif !B)$ 是公式。}{}

  \tagitem{prvIff}{若 $!A$ 和 $!B$ 是公式，则 $(!A \liff !B)$ 是公式。}{}

  \item 若 $!A$ 是公式且 $a \in G$，则 $\Knows_a !A$ 是公式。

  \item 若 $!A$ 和 $!B$ 是公式，则 $[!A] !B$ 是公式。

  \tagitem{limitClause}{除此之外没有其他公式。}{}
\end{enumerate}
\end{defn}

公式 $[!A] !B$ 的预期读法是“在 $!A$ 被真实宣告之后，$!B$ 成立”。有时在公开宣告的语境下讨论公共知识也是有用的，因此该语言也可包含公共知识算子~$\CKnows_G !A$。

\end{document}